\chapter{2004年天津卷（理）}
\section{单选题}
\begin{problem}
\(\i\) 是虚数单位，\(\frac{(-1 + \i)(2 + \i)}{\i^3} =\)（\quad）
\choices
    {\(1 + \i\)}
    {\(-1 - \i\)}
    {\(1 + 3\i\)}
    {\(-1 - 3\i\)}
\end{problem}

\begin{answer}
D
\end{answer}

\begin{solution}
因为 \(\i^3=-\i\)，且
\[
(-1+\i)(2+\i)=-2-\i+2\i+\i^2=-3+\i
\]
所以
\[
\frac{(-1+\i)(2+\i)}{\i^3}=\frac{-3+\i}{-\i}=(-3+\i)\i=-1-3\i
\]
故选 D.
\end{solution}

\begin{problem}
不等式 \(\frac{x - 1}{x} \geqslant 2\) 的解集为（\quad）

\choices
    {\([-1,0)\)}
    {\([-1, +\infty)\)}
    {\((-\infty, -1]\)}
    {\((-\infty, -1] \cup (0, +\infty)\)}
\end{problem}

\begin{answer}
A
\end{answer}

\begin{solution}
由 \(x\ne0\)，原不等式等价于 \(\frac{x-1}{x}-2=-\frac{x+1}{x}\geqslant0\)，即 \(\frac{x+1}{x}\leqslant0\). 临界点为 \(-1\) 和 \(0\). 当 \(-1\leqslant x<0\) 时该不等式成立，在其余区间不成立，因此解集为 \([-1,0)\).
故选 A.
\end{solution}

\begin{problem}
若平面向量 \(\bs{b}\) 与向量 \(\bs{a} = (1, -2)\) 的夹角是 \(180^{\circ}\)，且 \(\left|\bs{b}\right| = 3\sqrt{5}\)，则 \(\bs{b} =\)（\quad）

\choices
    {\((-3,6)\)}
    {\((3, -6)\)}
    {\((6, -3)\)}
    {\((-6,3)\)}
\end{problem}

\begin{answer}
A
\end{answer}

\begin{solution}
向量 \(\bs{b}\) 与 \(\bs{a}\) 的夹角为 \(180^{\circ}\)，所以存在 \(t<0\)，使得 \(\bs{b}=t\bs{a}\). 又
\(|\bs{a}|=\sqrt{1^2+(-2)^2}=\sqrt5\). 因此 \(|\bs{b}|=|t|\,|\bs{a}|=3\sqrt5\)，从而 \(|t|=3\). 由 \(t<0\) 得 \(t=-3\)，故 \(\bs{b}=-3(1,-2)=(-3,6)\).
故选 A.
\end{solution}

\begin{problem}
设 \(P\) 是双曲线 \(\frac{x^2}{a^2} - \frac{y^2}{9} = 1\) 上一点，双曲线的一条渐近线方程为 \(3x - 2y = 0\)，\(F_1\)，\(F_2\) 分别是双曲线的左，右焦点，若 \(\left|PF_1\right| = 3\)，则 \(\left|PF_2\right| =\)（\quad）

\choices
    {\(1\) 或 \(5\)}
    {\(6\)}
    {\(7\)}
    {\(9\)}
\end{problem}

\begin{answer}
C
\end{answer}

\begin{solution}
双曲线的渐近线斜率的绝对值为 \(\frac{3}{a}\)，而已知渐近线 \(3x-2y=0\) 的斜率为 \(\frac32\)，所以 \(\frac{3}{a}=\frac32\)，从而 \(a=2\). 因此双曲线的焦半距满足 \(c^2=a^2+9=4+9=13\).
双曲线上任一点到两焦点距离之差的绝对值为 \(2a=4\). 设 \(|PF_2|=d\)，由 \(|PF_1|=3\) 得 \(|3-d|=4\).
由于距离为正，只能得到 \(d=7\). 故选 C.
\end{solution}

\begin{problem}
若函数 \(f(x) = \log_a x\)（\(0 < a < 1\)）在区间 \([a, 2a]\) 上的最大值是最小值的 \(3\) 倍，则 \(a =\)（\quad）

\choices
    {\(\frac{\sqrt{2}}{4}\)}
    {\(\frac{\sqrt{2}}{2}\)}
    {\(\frac{1}{4}\)}
    {\(\frac{1}{2}\)}
\end{problem}

\begin{answer}
A
\end{answer}

\begin{solution}
因为 \(0<a<1\)，函数 \(f(x)=\log_a x\) 在正数范围内是减函数，所以在 \([a,2a]\) 上
\(f_{\max}=f(a)=1\)，\(f_{\min}=f(2a)=1+\log_a2\).
由最大值是最小值的 \(3\) 倍，得 \(1=3(1+\log_a2)\)，所以 \(\log_a2=-\frac23\). 于是 \(a^{-2/3}=2\)，从而 \(a=2^{-3/2}=\frac{\sqrt2}{4}\).
故选 A.
\end{solution}

\begin{problem}
如图，在棱长为 \(2\) 的正方体 \(ABCD - A_{1}B_{1}C_{1}D_{1}\) 中，\(O\) 是底面 \(ABCD\) 的中心，\(E\)，\(F\) 分别是 \(CC_{1}\)，\(AD\) 的中点. 那么异面直线 \(OE\) 和 \(FD_{1}\) 所成的角的余弦值等于（\quad）

\bitmapfigure[max width=0.34\linewidth]{img/2004/tianjin_science/q6_fig1.png}

\choices
    {\(\frac{\sqrt{10}}{5}\)}
    {\(\frac{\sqrt{15}}{5}\)}
    {\(\frac{4}{5}\)}
    {\(\frac{2}{5}\)}
\end{problem}

\begin{answer}
B
\end{answer}

\begin{solution}
建立如图所示直角坐标系，取
\(A=(0,0,0)\)、\(B=(2,0,0)\)、\(C=(2,2,0)\)、\(D=(0,2,0)\)，以及
\(A_1=(0,0,2)\)、\(B_1=(2,0,2)\)、\(C_1=(2,2,2)\)、\(D_1=(0,2,2)\). 则
\(O=(1,1,0)\)、\(E=(2,2,1)\)、\(F=(0,1,0)\).
可取异面直线 \(OE\) 和 \(FD_1\) 的方向向量分别为
\(\overrightarrow{OE}=(1,1,1)\) 和 \(\overrightarrow{FD_1}=(0,1,2)\). 因此
\[
\cos\theta
=\frac{\overrightarrow{OE}\cdot\overrightarrow{FD_1}}{|\overrightarrow{OE}|\,|\overrightarrow{FD_1}|}
=\frac{3}{\sqrt3\sqrt5}=\frac{\sqrt{15}}5
\]
故选 B.
\end{solution}

\begin{problem}
若 \(P(2, -1)\) 为圆 \((x - 1)^2 + y^2 = 25\) 的弦 \(AB\) 的中点，则直线 \(AB\) 的方程是（\quad）

\choices
    {\(x - y - 3 = 0\)}
    {\(2x + y - 3 = 0\)}
    {\(x + y - 1 = 0\)}
    {\(2x - y - 5 = 0\)}
\end{problem}

\begin{answer}
A
\end{answer}

\begin{solution}
圆心为 \(C_0=(1,0)\)，弦的中点为 \(P=(2,-1)\). 圆心与弦中点的连线垂直于弦，因此 \(\overrightarrow{C_0P}=(1,-1)\) 垂直于直线 \(AB\)，所以直线 \(AB\) 的方向向量可取 \((1,1)\). 直线过 \(P\)，故 \(y+1=x-2\)，即 \(x-y-3=0\).
故选 A.
\end{solution}

\begin{problem}
已知数列 \(\{a_{n}\}\)，那么“对任意的 \(n \in \N^{*}\)，点 \(P_{n}(n, a_{n})\) 都在直线 \(y = 2x + 1\) 上”是“\(\{a_{n}\}\) 为等差数列”的（\quad）

\choices
    {必要而不充分条件}
    {充分而不必要条件}
    {充要条件}
    {既不充分也不必要条件}
\end{problem}

\begin{answer}
B
\end{answer}

\begin{solution}
若对任意 \(n\in\N^*\)，点 \(P_n(n,a_n)\) 在直线 \(y=2x+1\) 上，则 \(a_n=2n+1\). 于是 \(a_{n+1}-a_n=\bigl(2(n+1)+1\bigr)-(2n+1)=2\)，
所以 \(\{a_n\}\) 是等差数列，故该条件是充分条件.

反过来，等差数列不一定满足 \(a_n=2n+1\). 例如取 \(a_n=0\)，它是等差数列，但点 \((n,0)\) 不在直线 \(y=2x+1\) 上. 因而该条件不是必要条件.
故选 B.
\end{solution}

\begin{problem}
函数 \(y = 2\sin \left(\frac{\pi}{6} - 2x\right)\)（\(x \in [0, \pi]\)）为增函数的区间是（\quad）
\choices
    {\(\left[0, \frac{\pi}{3}\right]\)}
    {\(\left[\frac{\pi}{12}, \frac{7\pi}{12}\right]\)}
    {\(\left[\frac{\pi}{3}, \frac{5\pi}{6}\right]\)}
    {\(\left[\frac{5\pi}{6}, \pi\right]\)}
\end{problem}

\begin{answer}
C
\end{answer}

\begin{solution}
函数的导数为
 \(y'=-4\cos\left(\frac{\pi}{6}-2x\right)\). 在 \(x\in[0,\pi]\) 内，函数增大的条件为 \(\cos\left(\frac{\pi}{6}-2x\right)\leqslant0\). 由于 \(\frac{\pi}{6}-2x\in[-\frac{11\pi}{6},\frac{\pi}{6}]\)，相关的余弦非正区间为 \(-\frac{3\pi}{2}\leqslant\frac{\pi}{6}-2x\leqslant-\frac{\pi}{2}\)，解得 \(\frac{\pi}{3}\leqslant x\leqslant\frac{5\pi}{6}\).
故选 C.
\end{solution}

\begin{problem}
如图，在长方体 \(ABCD - A_{1}B_{1}C_{1}D_{1}\) 中，\(AB = 6\)，\(AD = 4\)，\(AA_{1} = 3\). 分别过 \(BC\)，\(A_{1}D_{1}\) 的两个平行截面将长方体分成三部分，其体积分别记为 \(V_{1} = V_{AEA_{1}-DFD_{1}}\)，\(V_{2} = V_{EBE_{1}A_{1}-FCF_{1}D_{1}}\)，\(V_{3} = V_{B_{1}E_{1}B-C_{1}F_{1}C}\). 若 \(V_{1}: V_{2}: V_{3} = 1: 4: 1\)，则截面 \(A_{1}EFD_{1}\) 的面积为（\quad）

\bitmapfigure[max width=0.42\linewidth]{img/2004/tianjin_science/q10_fig1.png}

\choices
    {\(4\sqrt{10}\)}
    {\(8\sqrt{3}\)}
    {\(4\sqrt{13}\)}
    {\(16\)}
\end{problem}

\begin{answer}
C
\end{answer}

\begin{solution}
建立直角坐标系，取
\(A=(0,0,0)\)、\(B=(6,0,0)\)、\(D=(0,4,0)\)、\(C=(6,4,0)\)，以及 \(A_1=(0,0,3)\)、\(D_1=(0,4,3)\). 设截面 \(A_1EFD_1\) 与底面的交线满足 \(AE=DF=u\). 由于该截面经过 \(A_1D_1\)，在 \(xOz\) 平面内对应的截线是从 \(A_1\) 到 \(E\) 的线段；另一个截面与它平行且经过 \(BC\)，所以其在 \(xOz\) 平面内的截线从 \(B\) 到 \(E_1\)，并满足 \(B_1E_1=u\).

左侧部分在 \(xOz\) 平面上的截面是直角三角形，其两条直角边长分别为 \(u\) 和 \(3\)，并沿 \(y\) 轴方向延伸长度 \(4\)，所以
\[
V_1=\frac12\cdot u\cdot3\cdot4=6u
\]
右侧部分同理有 \(V_3=6u\). 长方体总体积为 \(6\cdot4\cdot3=72\)，故 \(V_2=72-V_1-V_3=72-12u\). 由体积比 \(V_1:V_2:V_3=1:4:1\)，得 \(72-12u=4(6u)\)，所以 \(u=2\).
在截面内，\(A_1D_1=4\)，且 \(\overrightarrow{A_1E}=(2,0,-3)\)、\(\overrightarrow{A_1D_1}=(0,4,0)\).
所以截面面积为
\[
S_{A_1EFD_1}=|\overrightarrow{A_1E}\times\overrightarrow{A_1D_1}|
=4\sqrt{2^2+3^2}=4\sqrt{13}
\]
故选 C.
\end{solution}

\begin{problem}
函数 \(y = 3^{x^2 - 1}\)（\(-1 \leqslant x < 0\)）的反函数是（\quad）

\choices
    {\(y = \sqrt{1 + \log_3 x}\)（\(x \geqslant \frac{1}{3}\)）}
    {\(y = -\sqrt{1 + \log_3 x}\)（\(x \geqslant \frac{1}{3}\)）}
    {\(y = \sqrt{1 + \log_3 x}\)（\(\frac{1}{3} < x \leqslant 1\)）}
    {\(y = -\sqrt{1 + \log_3 x}\)（\(\frac{1}{3} < x \leqslant 1\)）}
\end{problem}

\begin{answer}
D
\end{answer}

\begin{solution}
当 \(-1\leqslant x<0\) 时，\(x^2-1\) 的取值范围为 \((-1,0]\)，所以
\(y=3^{x^2-1}\in\left(\frac13,1\right]\). 由 \(\log_3y=x^2-1\) 得 \(x^2=1+\log_3y\). 原函数的定义域满足 \(x<0\)，故反函数中应取负平方根：\(x=-\sqrt{1+\log_3y}\). 交换自变量和因变量，反函数为 \(y=-\sqrt{1+\log_3x}\)，定义域为 \(\frac13<x\leqslant1\).
故选 D.
\end{solution}

\begin{problem}
定义在 \(\R\) 上的函数 \(f(x)\) 既是偶函数又是周期函数，若 \(f(x)\) 的最小正周期是 \(\pi\)，且当 \(x \in \left[0, \frac{\pi}{2}\right]\) 时，\(f(x) = \sin x\)，则 \(f\left(\frac{5\pi}{3}\right)\) 的值为（\quad）

\choices
    {\(-\frac{1}{2}\)}
    {\(\frac{1}{2}\)}
    {\(-\frac{\sqrt{3}}{2}\)}
    {\(\frac{\sqrt{3}}{2}\)}
\end{problem}

\begin{answer}
D
\end{answer}

\begin{solution}
由周期性，且周期为 \(\pi\)，有 \(f\left(\frac{5\pi}{3}\right)=f\left(\frac{2\pi}{3}\right)\). 又因为 \(f\) 是偶函数，\(f\left(\frac{2\pi}{3}\right)=f\left(-\frac{2\pi}{3}\right)=f\left(\frac{\pi}{3}\right)\)，其中最后一个等式再次利用了周期性. 由于 \(\frac{\pi}{3}\in[0,\frac{\pi}{2}]\)，所以 \(f\left(\frac{5\pi}{3}\right)=f\left(\frac{\pi}{3}\right)=\sin\frac{\pi}{3}=\frac{\sqrt3}{2}\).
故选 D.
\end{solution}

\section{填空题}
\begin{problem}
某工厂生产 \(A\)，\(B\)，\(C\) 三种不同型号的产品，产品数量之比依次为 \(2:3:5\)，现用分层抽样方法抽出一个容量为 \(n\) 的样本，样本中 \(A\) 种型号产品有 \(16\) 件. 那么此样本的容量 \(n =\)\fillinblank{}.
\end{problem}

\begin{answer}
80
\end{answer}

\begin{solution}
三种产品的总体比例为 \(2:3:5\)，其中 \(A\) 型号占总数的 \(\frac{2}{2+3+5}=\frac15\). 分层抽样保持各层比例不变，因此 \(\frac{16}{n}=\frac15\)，解得 \(n=80\).
\end{solution}

\begin{problem}
如果过两点 \(A(a,0)\) 和 \(B(0,a)\) 的直线与抛物线 \(y = x^{2} - 2x - 3\) 没有交点，那么实数 \(a\) 的取值范围是\fillinblank{}.
\end{problem}

\begin{answer}
\(a < -\frac{13}{4}\)
\end{answer}

\begin{solution}
当 \(a\ne0\) 时，过 \(A(a,0)\) 和 \(B(0,a)\) 的直线方程为
 \(x+y=a\)，即 \(y=a-x\). 代入抛物线方程，得 \(a-x=x^2-2x-3\)，即 \(x^2-x-(a+3)=0\).
直线与抛物线没有交点，等价于上述关于 \(x\) 的方程没有实根，因此 \(\Delta=(-1)^2-4\cdot1\cdot(-a-3)=4a+13\). 故没有交点的条件是 \(4a+13<0\)，即 \(a<-\frac{13}{4}\).
\end{solution}

\begin{problem}
若 \( (1-2x)^{2004}=a_{0}+a_{1}x+a_{2}x^{2}+\ldots+a_{2004}x^{2004}\)（\(x\in\R\)），则 \( (a_{0}+a_{1})+(a_{0}+a_{2})+(a_{0}+a_{3})+\cdots+(a_{0}+a_{2004})= \)\fillinblank{}.
\end{problem}

\begin{answer}
2004
\end{answer}

\begin{solution}
记
\[
G(x)=(1-2x)^{2004}=a_0+a_1x+a_2x^2+\cdots+a_{2004}x^{2004}
\]
由 \(G(0)=1\) 得 \(a_0=1\)，由 \(G(1)=1\) 得
\[
a_0+a_1+a_2+\cdots+a_{2004}=1
\]
题中所求的和为
\[
S=(a_0+a_1)+(a_0+a_2)+\cdots+(a_0+a_{2004})
\]
又因为
\[
a_1+a_2+\cdots+a_{2004}=G(1)-a_0
\]
所以
\[
S=2004a_0+(G(1)-a_0)=2003a_0+G(1)=2004
\]
\end{solution}

\begin{problem}
从 \(1\)，\(3\)，\(5\)，\(7\) 中任取 \(2\) 个数字，从 \(0\)，\(2\)，\(4\)，\(6\)，\(8\) 中任取 \(2\) 个数字，组成没有重复数字的四位数，其中能被 \(5\) 整除的四位数共有\fillinblank{}.
\end{problem}

\begin{answer}
300
\end{answer}

\begin{solution}
末位是 \(0\) 时，先从 \(1\)，\(3\)，\(5\)，\(7\) 中选出 \(2\) 个奇数，再从 \(2\)，\(4\)，\(6\)，\(8\) 中选出另一个偶数，最后排列其余三个非零数字，得到
\[
\mathrm C_4^2\cdot4\cdot3!=144
\]

末位是 \(5\) 时，还需从 \(1\)，\(3\)，\(7\) 中选出一个奇数，并从偶数集合中选出两个偶数. 若两个偶数中含 \(0\)，有 \(4\) 种选法，剩余三个数字的排列中需排除首位为 \(0\) 的情形，故每种奇数选择有 \(4(3!-2!)\) 个数；若不含 \(0\)，有 \(\mathrm C_4^2=6\) 种选法，三个非零数字可任意排列，故每种奇数选择有 \(6\cdot3!\) 个数. 因而这一类共有
\[
3\bigl[4(3!-2!)+\mathrm C_4^2\cdot3!\bigr]=156
\]
所以所求四位数共有
\[
144+156=300
\]
\end{solution}

\section{解答题}
\begin{problem}
已知 \(\tan \left(\frac{\pi}{4} + \alpha\right) = \frac{1}{2}\).

\begin{enumerate}
\item 求 \(\tan \alpha\) 的值；
\item 求 \(\frac{\sin 2\alpha - \cos^2\alpha}{1 + \cos 2\alpha}\) 的值.
\end{enumerate}
\end{problem}

\begin{solution}
\begin{enumerate}
\item 令 \(t=\tan\alpha\). 由正切加法公式，\(\frac{1+t}{1-t}=\frac12\)，所以 \(2+2t=1-t\)，解得 \(t=-\frac13\). 因而 \(\tan\alpha=-\frac13\).

\item 因为 \(1+\cos2\alpha=2\cos^2\alpha\)，所以
\[
\frac{\sin2\alpha-\cos^2\alpha}{1+\cos2\alpha}
=\frac{2\sin\alpha\cos\alpha-\cos^2\alpha}{2\cos^2\alpha}
=\tan\alpha-\frac12
=-\frac13-\frac12=-\frac56
\]
\end{enumerate}
\end{solution}

\begin{problem}
从 \(4\) 名男生和 \(2\) 名女生中任选 \(3\) 人参加演讲比赛，设随机变量 \(\xi\) 表示所选 \(3\) 人中女生的人数.

\begin{enumerate}
\item 求 \(\xi\) 的分布列；
\item 求 \(\xi\) 的数学期望；
\item 求“所选 \(3\) 人中女生人数 \(\xi \leqslant 1\)”的概率.
\end{enumerate}
\end{problem}

\begin{solution}
\begin{enumerate}
\item 从 \(6\) 人中任选 \(3\) 人，共有 \(\mathrm C_6^3=20\) 种等可能结果. 随机变量 \(\xi\) 的可能取值为 \(0\)，\(1\)，\(2\). 当 \(\xi=0\) 时，从 \(4\) 名男生中任选 \(3\) 人；当 \(\xi=1\) 时，从 \(2\) 名女生中任选 \(1\) 人、从 \(4\) 名男生中任选 \(2\) 人；当 \(\xi=2\) 时，选中 \(2\) 名女生和 \(1\) 名男生. 因此 \(P(\xi=0)=\frac{\mathrm C_4^3}{\mathrm C_6^3}=\frac15\)，\(P(\xi=1)=\frac{\mathrm C_2^1\mathrm C_4^2}{\mathrm C_6^3}=\frac35\)，\(P(\xi=2)=\frac{\mathrm C_2^2\mathrm C_4^1}{\mathrm C_6^3}=\frac15\). 分布列为
\[
\begin{tabular}{c|ccc}
\(\xi\) & \(0\) & \(1\) & \(2\) \\
\hline
\(P(\xi)\) & \(\frac15\) & \(\frac35\) & \(\frac15\)
\end{tabular}
\]

\item \(E(\xi)=0\cdot\frac15+1\cdot\frac35+2\cdot\frac15=1\).

\item 所求概率为 \(P(\xi\leqslant1)=P(\xi=0)+P(\xi=1)=\frac15+\frac35=\frac45\).
\end{enumerate}
\end{solution}

\begin{problem}
如图，在四棱锥 \(P - ABCD\) 中，底面 \(ABCD\) 是正方形，侧棱 \(PD \perp\) 底面 \(ABCD\)，\(PD = DC\)，\(E\) 是 \(PC\) 的中点，作 \(EF \perp PB\) 交 \(PB\) 于点 \(F\).

\begin{enumerate}
\item 证明：\(PA \parallel\) 平面 \(EDB\)；
\item 证明：\(PB \perp\) 平面 \(EFD\)；
\item 求二面角 \(C - PB - D\) 的大小.
\end{enumerate}

\begin{center}
\bitmapfigure[width=13.3cm]{img/2004/tianjin_science/q19_fig1.png}
\end{center}
\end{problem}

\begin{solution}
\begin{enumerate}
\item 设正方形边长为 \(s\)，建立直角坐标系，取
\(D=(0,0,0)\)，\(C=(s,0,0)\)，\(A=(0,s,0)\)，\(B=(s,s,0)\)，\(P=(0,0,s)\).
由 \(E\) 是 \(PC\) 的中点，得 \(E=(\frac{s}{2},0,\frac{s}{2})\). 因而
\(\overrightarrow{PA}=(0,s,-s)\)，\(\overrightarrow{DB}=(s,s,0)\)，\(\overrightarrow{DE}=(\frac{s}{2},0,\frac{s}{2})\)，并且 \(\overrightarrow{PA}=\overrightarrow{DB}-2\overrightarrow{DE}\).

由于 \(\overrightarrow{DB}\) 和 \(\overrightarrow{DE}\) 是平面 \(EDB\) 内两个相交且不共线的方向向量，且 \(\overrightarrow{PA}\) 可由它们线性表示，直线 \(PA\) 的方向平行于平面 \(EDB\). 又平面 \(EDB\) 的方程为 \(-x+y+z=0\)，而点 \(P=(0,0,s)\) 不在此平面内，所以 \(PA\parallel\) 平面 \(EDB\).

\item \(\overrightarrow{PB}=(s,s,-s)\)，\(\overrightarrow{DE}=(\frac{s}{2},0,\frac{s}{2})\)，所以 \(\overrightarrow{PB}\cdot\overrightarrow{DE}=\frac{s^2}{2}-\frac{s^2}{2}=0\)，即 \(PB\perp ED\). 题设又给出 \(PB\perp EF\)，而 \(ED\) 与 \(EF\) 是平面 \(EFD\) 内相交的两条直线，因此 \(PB\perp\) 平面 \(EFD\).

\item 求二面角 \(C-PB-D\). 平面 \(CPB\) 的一个法向量为
\[
\overrightarrow{PC}\times\overrightarrow{PB}=s^2(1,0,1)
\]
平面 \(DPB\) 的一个法向量为
\[
\overrightarrow{PD}\times\overrightarrow{PB}=s^2(1,-1,0)
\]
设二面角的较小角为 \(\theta\)，则
\(\cos\theta=\frac{(1,0,1)\cdot(1,-1,0)}{\sqrt2\sqrt2}=\frac12\)，所以 \(\theta=60^{\circ}\).
\end{enumerate}
\end{solution}

\begin{problem}
已知函数 \(f(x) = ax^3 + bx^2 - 3x\) 在 \(x = \pm 1\) 处取得极值.

\begin{enumerate}
\item 讨论 \(f(1)\) 和 \(f(-1)\) 是函数 \(f(x)\) 的极大值还是极小值；
\item 过点 \(A(0,16)\) 作曲线 \(y = f(x)\) 的切线. 求此切线方程.
\end{enumerate}
\end{problem}

\begin{solution}
\begin{enumerate}
\item 由 \(f(x)=ax^3+bx^2-3x\)，得 \(f'(x)=3ax^2+2bx-3\). 在 \(x=1\) 和 \(x=-1\) 处取得极值，所以 \(3a+2b-3=0\)，\(3a-2b-3=0\). 两式相加得 \(6a-6=0\)，于是 \(a=1\)；再代回得 \(b=0\). 因而 \(f(x)=x^3-3x\)，\(f'(x)=3(x^2-1)\).

当 \(x<-1\) 或 \(x>1\) 时 \(f'(x)>0\)，当 \(-1<x<1\) 时 \(f'(x)<0\). 所以函数在 \(x=-1\) 处先增后减，取得极大值 \(f(-1)=2\)；在 \(x=1\) 处先减后增，取得极小值 \(f(1)=-2\).

\item 设切点横坐标为 \(t\). 切线为 \(y=f(t)+f'(t)(x-t)\)，其在 \(y\) 轴上的截距为
\[
f(t)-tf'(t)=(t^3-3t)-t(3t^2-3)=-2t^3
\]
切线过 \(A(0,16)\)，故 \(-2t^3=16\)，从而 \(t=-2\). 此时斜率 \(f'(-2)=3((-2)^2-1)=9\)，所以所求切线方程为 \(y=9x+16\).
\end{enumerate}
\end{solution}

\begin{problem}
已知定义在 \(\R\) 上的函数 \(f(x)\) 和数列 \(\{a_n\}\) 满足下列条件：

\begin{circlelist}
\item \(a_1 = a\)，\(a_n = f(a_{n-1})\)（\(n = 2\)，\(3\)，\(4\)，\(\cdots\)）；
\item \(a_2 \neq a_1\)，\(f(a_n) - f(a_{n-1}) = k(a_n - a_{n-1})\)（\(n = 2\)，\(3\)，\(4\)，\(\cdots\)）；
\item 其中 \(a\) 为常数，\(k\) 为非零常数.
\end{circlelist}

\begin{enumerate}
\item 令 \(b_{n} = a_{n+1} - a_{n}\)（\(n \in \N^{*}\)），证明数列 \(\{b_{n}\}\) 是等比数列；
\item 求数列 \(\{a_{n}\}\) 的通项公式；
\item 当 \(|k| < 1\) 时，求 \(\displaystyle\lim_{n \to \infty} a_n\).
\end{enumerate}
\end{problem}

\begin{solution}
\begin{enumerate}
\item 由递推关系 \(a_{n+1}=f(a_n)\) 和题设条件，对于 \(n\geqslant2\)，有
\[
b_n=a_{n+1}-a_n=f(a_n)-f(a_{n-1})=k(a_n-a_{n-1})=kb_{n-1}
\]
又因 \(b_1=a_2-a_1\ne0\)，且 \(k\ne0\)，所以 \(\{b_n\}\) 是首项为 \(b_1=a_2-a\)、公比为 \(k\) 的等比数列，即 \(b_n=(a_2-a)k^{n-1}\).

\item 由 \(a_n=a_1+b_1+b_2+\cdots+b_{n-1}\)，当 \(k\ne1\) 时，
\[
a_n=a+(a_2-a)(1+k+\cdots+k^{n-2})
=a+(a_2-a)\frac{1-k^{n-1}}{1-k}
\]
当 \(k=1\) 时，\(b_n=a_2-a\)，从而 \(a_n=a+(n-1)(a_2-a)\).

\item 当 \(|k|<1\) 时必有 \(k\ne1\)，且 \(k^{n-1}\to0\)，所以
\[
\lim_{n\to\infty}a_n
=a+\frac{a_2-a}{1-k}
=\frac{a_2-ak}{1-k}
\]
\end{enumerate}
\end{solution}

\begin{problem}
椭圆的中心是原点 \(O\)，它的短轴长为 \(2\sqrt{2}\). 相应于焦点 \(F(c,0)\)（\(c > 0\)）的准线 \(l\) 与 \(x\) 轴相交于点 \(A\)，\(|OF| = 2|FA|\). 过点 \(A\) 的直线与椭圆相交于 \(P\)，\(Q\) 两点.

\begin{enumerate}
\item 求椭圆的方程及离心率；
\item 若 \(\overrightarrow{OP} \cdot \overrightarrow{OQ} = 0\)，求直线 \(PQ\) 的方程；
\item 设 \(\overrightarrow{AP} = \lambda \overrightarrow{AQ}\)（\(\lambda > 1\)），过点 \(P\) 且平行于准线 \(l\) 的直线与椭圆相交于另一点 \(M\)，证明 \(\overrightarrow{FM} = -\lambda \overrightarrow{FQ}\).
\end{enumerate}
\end{problem}

\begin{solution}
\begin{enumerate}
\item 设椭圆的长半轴为 \(a\)，短半轴为 \(b\). 由短轴长为 \(2\sqrt2\)，得 \(b=\sqrt2\)，即 \(b^2=2\). 由于焦点为 \(F(c,0)\)，椭圆方程可设为
\(\frac{x^2}{a^2}+\frac{y^2}{2}=1\)，\(c^2=a^2-2\).
对应于右焦点 \(F(c,0)\) 的准线为 \(l:x=\frac{a^2}{c}\)，所以
\(A=\left(\frac{a^2}{c},0\right)\)，\(|FA|=\frac{a^2}{c}-c=\frac{a^2-c^2}{c}=\frac2c\).
由 \(|OF|=2|FA|\)，得 \(c=2\cdot\frac2c\). 因为 \(c>0\)，所以 \(c=2\). 于是 \(a^2=c^2+2=6\)，故椭圆方程为 \(\frac{x^2}{6}+\frac{y^2}{2}=1\)，离心率为 \(e=\frac ca=\frac2{\sqrt6}=\frac{\sqrt6}{3}\).

\item 因为准线与 \(x\) 轴交于 \(A=(3,0)\)，设直线 \(PQ\) 的斜率为 \(m\). 竖直直线 \(x=3\) 不与椭圆相交，所以不遗漏这种情形. 直线方程为 \(y=m(x-3)\)，代入椭圆方程得
\[
(1+3m^2)x^2-18m^2x+27m^2-6=0
\]
设两个交点的横坐标为 \(x_P\)，\(x_Q\)，则
\(x_P+x_Q=\frac{18m^2}{1+3m^2}\)，\(x_Px_Q=\frac{27m^2-6}{1+3m^2}\).
又因为 \(y_P=m(x_P-3)\)，\(y_Q=m(x_Q-3)\)，有
\[
(x_P-3)(x_Q-3)=x_Px_Q-3(x_P+x_Q)+9=\frac3{1+3m^2}
\]
因此
\[
\overrightarrow{OP}\cdot\overrightarrow{OQ}
=x_Px_Q+y_Py_Q
=x_Px_Q+m^2(x_P-3)(x_Q-3)
=\frac{6(5m^2-1)}{1+3m^2}
\]
由条件 \(\overrightarrow{OP}\cdot\overrightarrow{OQ}=0\)，得 \(m^2=\frac15\). 所以直线 \(PQ\) 为 \(x-\sqrt5y-3=0\) 或 \(x+\sqrt5y-3=0\).

\item 设 \(\overrightarrow{AQ}=(u,v)\). 由 \(\overrightarrow{AP}=\lambda\overrightarrow{AQ}\)，有 \(Q=(3+u,v)\)，\(P=(3+\lambda u,\lambda v)\). 以 \(A\) 为起点，将该直线上的点写成 \(X(t)=(3+tu,tv)\). 交点参数 \(t=1\)，\(\lambda\) 代入椭圆方程，得到
\[
(u^2+3v^2)t^2+6ut+3=0
\]
由韦达定理，因其两个根为 \(1\) 和 \(\lambda\)，所以
\(u^2+3v^2=\frac3\lambda\)，\(1+\lambda=-\frac{6u}{u^2+3v^2}=-2\lambda u\).
即 \(1+\lambda u=-\lambda(1+u)\). 过 \(P\) 且平行于准线 \(l:x=3\) 的直线为竖直线，椭圆关于 \(x\) 轴对称，因此它与椭圆的另一交点为 \(M=(3+\lambda u,-\lambda v)\). 由于 \(F=(2,0)\)，于是 \(\overrightarrow{FM}=(1+\lambda u,-\lambda v)\)，而 \(-\lambda\overrightarrow{FQ}=-\lambda(1+u,v)=(-\lambda(1+u),-\lambda v)\). 利用 \(1+\lambda u=-\lambda(1+u)\)，可得 \(\overrightarrow{FM}=-\lambda\overrightarrow{FQ}\).
\end{enumerate}
\end{solution}
